\documentclass[11pt]{article}
\usepackage[T1]{fontenc}
\usepackage[utf8]{inputenc}
\usepackage{lmodern,amsmath,amssymb,amsthm,mathtools}
\usepackage[margin=1in]{geometry}
\usepackage{booktabs,array,longtable,microtype}
\usepackage[hidelinks,breaklinks]{hyperref}
\usepackage{xurl}
\hypersetup{pdftitle={Exact qutrit APPT purity and the unrestricted higher-dimensional asymptotic law},pdfauthor={Yongxian Zhang},pdfsubject={APPT purity: exact qutrit theorem, counterexample, and asymptotic law},pdfkeywords={absolute PPT, purity, spectra, asymptotics, exact certificates, Lean}}
\newtheorem{theorem}{Theorem}[section]
\newtheorem{lemma}[theorem]{Lemma}
\newtheorem{proposition}[theorem]{Proposition}
\newtheorem{corollary}[theorem]{Corollary}
\theoremstyle{remark}\newtheorem{remark}[theorem]{Remark}
\newcommand{\Tr}{\operatorname{Tr}}
\newcommand{\APPT}{\operatorname{APPT}}
\newcommand{\id}{\mathrm{id}}
\newcommand{\C}{\mathbb C}
\newcommand{\R}{\mathbb R}
\newcommand{\Q}{\mathbb Q}
\newcommand{\N}{\mathbb N}
\newcommand{\psd}{\succeq0}
\newcommand{\hs}{\mathrm{HS}}
\newcommand{\proofcommit}{\texttt{2fc3f25179cf1128c4bccbde65fd3a07c093ad6d}}
\title{Exact qutrit APPT purity and the unrestricted\protect\\higher-dimensional asymptotic law}
\author{Yongxian Zhang\thanks{School of Computer Science and Engineering, South China University of Technology, Guangzhou, China. Correspondence: \href{mailto:mxymmxym1@gmail.com}{mxymmxym1@gmail.com}. ORCID: \url{https://orcid.org/0009-0000-3864-3536}.}}
\date{10 October 2026}
\begin{document}
\maketitle
\begin{abstract}
We give a unified set of results on the purity of absolutely positive partial transpose (APPT) states. First, for every integer $n\geq3$, the exact maximum on $\C^3\otimes\C^n$ is $(3n+8)/(3n+2)^2$ for $3\leq n\leq8$ and $3/(8n)$ for $n\geq9$. This proves the qutrit specialization of the maximal-value and attainment conjecture of Ahiable, Kothakonda and Winter. The entire qutrit theorem, including the physical global-unitary quantifier and actual attainment, is formalized in Lean 4.34.1 and frozen in an immutable proof release. Second, the proposed general-dimensional formula is false: an explicit APPT state on $\C^{10}\otimes\C^{38}$ has purity strictly larger than both conjectured candidates, and an infinite family gives further violations. We then prove a single corrected leading law whenever $D=mn\to\infty$, uniformly over all $2\leq m\leq n$:\[ P_{\max}(m,n)-\frac1D\sim\frac{\max\{8,4+D/(m^2-1)\}}{D^2}.\] For $m\to\infty$ this becomes $D^2(P_{\max}-1/D)\sim\max\{8,4+n/m\}$, while fixed $m$ gives $D(P_{\max}-1/D)\to1/(m^2-1)$. The higher-dimensional arguments are written analytic proofs supported by exact rational checkers; they are not included in the qutrit Lean certificate. For every $m\geq3$ and $n\geq m^3-m-2$, we also determine the exact finite-dimensional maximum and all maximizing spectra. Exact finite-dimensional maximizers outside that rectangular range and the absolute-separability maximum in general dimensions remain open.
\end{abstract}
\noindent\textbf{Keywords.} Absolute PPT; purity; spectra; quantum states; exact certificates; Lean formalization.

\section{Introduction and main result}
Fix a tensor product $\C^3\otimes\C^n$ and let $\Gamma=T\otimes\id$ denote partial transposition on its first factor. A density matrix $\rho$ is \emph{absolutely PPT} if
\begin{equation}\label{eq:appt}
 (U\rho U^*)^\Gamma\succeq0\qquad\text{for every }U\in\mathcal U(3n).
\end{equation}
Unlike the PPT property in a fixed basis, this condition depends only on the eigenvalues. Purity, $\Tr(\rho^2)$, measures the squared Hilbert--Schmidt norm. Determining its maximum is equivalent to determining the smallest ball centered at the maximally mixed state that contains the APPT set.

Hildebrand characterized the APPT property by spectral linear matrix inequalities \cite{Hildebrand}. Dung and Khoi investigated maximum purity and proposed numerical formulas for qubit--qudit and qutrit--qudit systems \cite{DungKhoi}. Their qutrit formula is not the formula proved here. Tran's later outer-polytope analysis gives a strict qutrit upper bound that disproves that earlier qutrit conjecture \cite{Tran}; we make no claim to the first disproof. Ahiable, Kothakonda and Winter computed the exact maximum over an inscribed APPT polytope and conjectured that it is also the maximum over the full APPT spectral set \cite[Theorem 6.3, Corollary 6.4 and Conjecture 6.7]{AKW}. Its specialization to local dimension three is the following theorem.

\begin{theorem}[Exact qutrit--qudit APPT purity]\label{thm:main}
For every integer $n\geq3$,
\begin{equation}\label{eq:maximum}
 \max_{\rho\in\APPT_{3,n}}\Tr(\rho^2)
 =p_n:={\begin{cases}
 \displaystyle\frac{3n+8}{(3n+2)^2},&3\leq n\leq8,\\[6pt]
 \displaystyle\frac{3}{8n},&n\geq9.
 \end{cases}}
\end{equation}
The first value is attained by $(I+2vv^*)/(3n+2)$ for any unit vector $v$. The second value is attained by $(I+P)/(4n)$ for any rank-$n$ orthogonal projection $P$.
\end{theorem}

Here $\APPT_{3,n}$ denotes the trace-one positive-semidefinite matrices satisfying \eqref{eq:appt}. The result includes both an upper bound for every such matrix and an actual attaining state. It confirms the value and provides attaining spectra from the conjecture of \cite{AKW}; it does not classify all states attaining equality. The two displayed witnesses have spectra
\begin{equation}\label{eq:witness-spectra}
 \frac{1}{3n+2}(3,1,\ldots,1),\qquad
 \frac{1}{4n}(\underbrace{2,\ldots,2}_{n},
                  \underbrace{1,\ldots,1}_{2n}),
\end{equation}
respectively. 
In particular, the asserted extremal values were predicted by previous work; the contribution here is their proof for the entire qutrit--qudit family and a complete physical-state formalization.

The same conjecture was stated for arbitrary local dimensions. The higher-dimensional result below shows that its proposed finite-dimensional formula cannot hold beyond the qutrit regime. We write $P_{\max}(m,n)$ for the maximum purity over all APPT density matrices on $\C^m\otimes\C^n$, where $2\leq m\leq n$. For the conjecture, put $D=mn$ and $t=\lceil(m-1)n/2\rceil$ and define
\begin{equation}\label{eq:general-candidate}
 M_{m,n}=\max\left\{\frac{D+8}{(D+2)^2},
 \frac{D(m-1)^2+4mt}{[D(m-1)+2t]^2}\right\}.
\end{equation}
The second part of this paper proves that $P_{\max}(m,n)>M_{m,n}$ for explicit dimensions and gives the asymptotic replacement for $M_{m,n}$ when $m$ grows. It also proves an exact finite-dimensional formula, including equality cases, in the sufficiently rectangular range $n\geq m^3-m-2$.

The qutrit proof is computer-assisted in a precise sense. Two polynomial nonnegativity lemmas are established by rational certificates supplied as finite data. One covers the six dimensions $9,12,\ldots,24$; the other covers arbitrary real endpoint populations with sum at least $18$. All remaining steps are analytic, linear-algebraic, or elementary algebraic. Sections~\ref{sec:finite} and~\ref{sec:uniform} state the exact identities and explain how to reconstruct and check every term. The Lean proof independently verifies these identities and the bridges to the physical theorem. Neither a numerical optimization result nor a solver success flag is a premise.

\section{Physical necessity and the nine selected eigenvalues}\label{sec:necessity}
Write $D=3n$ and order the eigenvalues of a density matrix as
\begin{equation}\label{eq:spectrum}
 \lambda_1\geq\lambda_2\geq\cdots\geq\lambda_D\geq0,
 \qquad \sum_{i=1}^D\lambda_i=1.
\end{equation}
For $n\geq3$, select the nine distinct positions
\begin{equation}\label{eq:outer}
 y=(\lambda_1,\lambda_2,\lambda_3,
       \lambda_{D-5},\lambda_{D-4},\lambda_{D-3},
       \lambda_{D-2},\lambda_{D-1},\lambda_D).
\end{equation}
Define the real symmetric matrices
\begin{equation}\label{eq:A}
 A(y)=\begin{pmatrix}
 2y_9&y_8-y_1&y_6-y_2\\
 y_8-y_1&2y_7&y_5-y_3\\
 y_6-y_2&y_5-y_3&2y_4
 \end{pmatrix},
\end{equation}
\begin{equation}\label{eq:B}
 B(y)=\begin{pmatrix}
 2y_9&y_8-y_1&y_7-y_2\\
 y_8-y_1&2y_6&y_5-y_3\\
 y_7-y_2&y_5-y_3&2y_4
 \end{pmatrix}.
\end{equation}
These are the familiar qutrit spectral tests, with the two matrices possibly interchanged relative to another labeling convention.

\begin{lemma}[Physical necessity]\label{lem:necessity}
If $\rho$ is APPT, then $A(y)\succeq0$ and $B(y)\succeq0$ for $y$ in \eqref{eq:outer}. More generally, the conclusion holds for any nine distinct eigenvalue positions assigned to the coordinates of $y$.
\end{lemma}
\begin{proof}
By a global unitary, first diagonalize $\rho$ and move the selected eigenvalues into a $3\times3$ product-basis corner. For $A$, assign those eigenvalues to the corner coordinates according to the array
\[
 \begin{pmatrix}9&8&6\\1&7&5\\2&3&4\end{pmatrix};
 \quad\text{for }B\text{ use }
 \begin{pmatrix}9&8&7\\1&6&5\\2&3&4\end{pmatrix}.
\]
An entry $j$ in an array means that the corresponding diagonal entry is $y_j$. The other eigenvalues occupy the complementary coordinates. On each of the three planes spanned by $e_i\otimes e_j,e_j\otimes e_i$ with $i<j$, apply the real unitary
\[
 \frac1{\sqrt2}\begin{pmatrix}1&1\\1&-1\end{pmatrix},
\]
and leave the diagonal corner coordinates and the orthogonal complement fixed. A diagonal pair with upper-coordinate value $s$ and lower-coordinate value $a$ acquires off-diagonal entries $(s-a)/2$. Partial transposition carries these entries to the matrix elements between $e_i\otimes e_i$ and $e_j\otimes e_j$. Thus the principal submatrix on the three diagonal corner vectors is exactly $A(y)/2$ or $B(y)/2$. Condition~\eqref{eq:appt} makes the whole partial transpose positive semidefinite; its principal submatrices are positive semidefinite as well.
\end{proof}

This argument only uses necessity. It does not assume that these matrices characterize APPT, and it does not need the converse direction of Hildebrand's criterion. In Lean, arbitrary injections of nine eigenvalue slots are extended to permutations of the full index type, the corner unitary is extended by the identity, and its entries are verified explicitly. Hermitian diagonalization then supplies the actual ordered spectrum, together with
\begin{equation}\label{eq:moments}
 \Tr(\rho)=\sum_i\lambda_i,\qquad
 \Tr(\rho^2)=\sum_i\lambda_i^2.
\end{equation}
These equations are proved from diagonalization, not supplied as spectral-interface hypotheses.

We shall use the following immediate PSD consequences. For $C=A(y)$ or $B(y)$, its determinant, every principal $2\times2$ minor and every real quadratic form $v^TCv$ are nonnegative. Also,
\begin{lemma}\label{lem:ratio}
If $y_1\geq\cdots\geq y_9\geq0$ and $A(y)\succeq0$, then $2y_4-y_3\geq0$.
\end{lemma}
\begin{proof}
Writing $\boldsymbol1=(1,1,1)^T$, we have
\[
 \tfrac12\boldsymbol1^TA(y)\boldsymbol1
 =\sum_{j=4}^9y_j-\sum_{j=1}^3y_j\geq0.
\]
Ordering gives $3y_3\leq\sum_{j=1}^3y_j\leq\sum_{j=4}^9y_j\leq6y_4$.
\end{proof}

\section{Six finite cubic certificates}\label{sec:finite}
For $D\in\{9,12,15,18,21,24\}$ introduce gaps
\[
 g_i=\lambda_i-\lambda_{i+1}\ (i<D),\qquad g_D=\lambda_D.
\]
Then $g_i\geq0$ and $\lambda_i=\sum_{j=i}^Dg_j$. Conversely every nonnegative gap vector defines an ordered nonnegative spectrum in this way. Set
\[
 T=\sum_{i=1}^D\lambda_i,\qquad S=\sum_{i=1}^D\lambda_i^2,
\]
and form $y$ using~\eqref{eq:outer}.

\begin{lemma}[Finite certificate identities]\label{lem:finite}
For each of the six dimensions above, the ancillary certificate specifies positive rational numbers $c_{D,\alpha}$ and explicitly named homogeneous cubic polynomials $F_{D,\alpha}$ such that the polynomial identity
\begin{equation}\label{eq:finite-cert}
 \bigl((D+8)T^2-(D+2)^2S\bigr)T
 =\sum_{\alpha}c_{D,\alpha}F_{D,\alpha}(g)
\end{equation}
holds over $\Q[g_1,\ldots,g_D]$. Every polynomial on the right is nonnegative whenever $g\geq0$, $A(y)\succeq0$ and $B(y)\succeq0$.
\end{lemma}
\begin{proof}[Certificate proof]
The allowable terms are products of nonnegative gaps with PSD determinants, principal minors or quadratic forms, and plain gap monomials. Explicitly, the data use terms of the forms
\begin{equation}\label{eq:finite-generators}
 \det C,\quad
 (C_{ii}C_{jj}-C_{ij}C_{ji})g_k,\quad
 (v^TCv)g_kg_\ell,\quad g_ig_jg_k,
 \qquad C\in\{A(y),B(y)\},
\end{equation}
with the rational vector $v$ specified in the term name. Each term is nonnegative under the hypotheses. The six rational coefficient lists are included as exact decimal-free strings in the ancillary JSON files. Table~\ref{tab:finite} specifies their sizes and the resulting bounds.

For verification, substitute the linear tail sums for $\lambda$, expand both sides of~\eqref{eq:finite-cert} over the rational numbers, and compare the coefficient of every monomial. The supplied checker performs precisely this operation using Python integers and \texttt{Fraction}; it also verifies strict positivity of every certificate weight. All six coefficient differences vanish. Equivalently, one may clear denominators by the least common multiple of the denominators and compare integer coefficient dictionaries. The Lean finite modules separately prove the same identities through small symbolic expansions and kernel-checked integer coefficient merges, and derive their nonnegativity from the PSD consequences described above. The data, not a numerical search procedure, constitute the finite part of the proof.
\end{proof}

\begin{table}[ht]
\centering
\begin{tabular}{rrrrl}
\toprule
$n$&$D$&Certificate terms&$(D+8)/(D+2)^2$&Data filename suffix\\
\midrule
3&9&166&$17/121$&\texttt{cubic\_certificate\_search.json}\\
4&12&363&$5/49$&\texttt{D12\_cubic\_certificate.json}\\
5&15&679&$23/289$&\texttt{D15\_fast\_certificate.json}\\
6&18&1139&$13/200$&\texttt{D18\_fast\_certificate.json}\\
7&21&1769&$29/529$&\texttt{D21\_fast\_certificate.json}\\
8&24&2596&$8/169$&\texttt{D24\_fast\_certificate.json}\\
\bottomrule
\end{tabular}
\caption{Finite exact certificates. Each filename has the prefix \texttt{qutrit\_}. Counts refer to positive weighted terms, not expanded monomials.}\label{tab:finite}
\end{table}

\begin{corollary}\label{cor:finite}
Every APPT density matrix on $\C^3\otimes\C^n$ with $3\leq n\leq8$ satisfies $\Tr(\rho^2)\leq(3n+8)/(3n+2)^2$.
\end{corollary}
\begin{proof}
Lemma~\ref{lem:necessity} supplies the two PSD hypotheses. Ordering supplies nonnegative gaps, and trace one gives $T=1$. Hence~\eqref{eq:finite-cert} yields $(D+8)-(D+2)^2S\geq0$. Apply~\eqref{eq:moments}.
\end{proof}

\section{A uniform bound by endpoint compression}\label{sec:uniform}
The number of middle eigenvalues grows with $n$, but their positions do not enter the two PSD matrices. The following elementary fact removes their individual values.
\begin{lemma}[Endpoint compression]\label{lem:compression}
Let $b\leq x_i\leq a$ for $1\leq i\leq M$, where $b\leq a$. There exist real $t,z\geq0$ such that
\begin{equation}\label{eq:compression}
 t+z=M,\quad \sum_i x_i=ta+zb,\quad
 \sum_i x_i^2\leq ta^2+zb^2.
\end{equation}
\end{lemma}
\begin{proof}
For $a>b$, choose
\[
 t=\frac{\sum_i x_i-Mb}{a-b},\qquad z=M-t.
\]
The interval constraints imply $0\leq t\leq M$. Since $(a-x_i)(x_i-b)\geq0$, we have $x_i^2\leq(a+b)x_i-ab$; sum this inequality to obtain~\eqref{eq:compression}. If $a=b$, take $t=M$ and $z=0$.
\end{proof}
The populations $t,z$ are real rather than necessarily integral. This is important: rounding them would not preserve the trace exactly.

For an arbitrary ordered $y\in\R^9_{\geq0}$ and $t,z\geq0$, write
\begin{align}
 T(y,t,z)&=\sum_{i=1}^9y_i+t y_3+z y_4,\label{eq:uniform-T}\\
 S(y,t,z)&=\sum_{i=1}^9y_i^2+t y_3^2+z y_4^2.\label{eq:uniform-S}
\end{align}
Take $g_i=y_i-y_{i+1}$ for $i<9$ and $g_9=y_9$, and abbreviate
\begin{equation}\label{eq:mixed}
 w=t+z-18,\quad h=t y_3-z y_4,\quad
 q=9T-4(9+t+z)(y_3+y_4).
\end{equation}

\begin{lemma}[Uniform rational certificate]\label{lem:uniform}
For every ordered nonnegative $y$, every $t,z\geq0$ with $t+z\geq18$, and $A(y),B(y)\succeq0$,
\begin{equation}\label{eq:uniform-nonneg}
 \bigl(9T(y,t,z)^2-8(9+t+z)S(y,t,z)\bigr)T(y,t,z)\geq0.
\end{equation}
In particular, if $T(y,t,z)=1$, then
\begin{equation}\label{eq:uniform-bound}
 S(y,t,z)\leq\frac{9}{8(9+t+z)}.
\end{equation}
\end{lemma}
\begin{proof}[Certificate proof]
The file \texttt{qutrit\_uniform\_parameter\_search.json} gives 1,635 positive rational weights $c_\alpha$ and explicitly named polynomials $F_\alpha$ satisfying
\begin{equation}\label{eq:uniform-cert}
 \bigl(9T^2-8(9+t+z)S\bigr)T
 =\sum_{\alpha=1}^{1635}c_\alpha F_\alpha(g,t,z).
\end{equation}
This is a polynomial identity in the nine gaps and the two real parameters. Its nonnegative generators consist of the PSD-based and gap terms in~\eqref{eq:finite-generators}, the additional terms $(2y_4-y_3)g_i g_j$, nonnegative monomials in the gaps and $t,z$, and terms of the form
\begin{equation}\label{eq:square-generators}
 h^2g_i\xi\quad\hbox{or}\quad (a q+b h)^2g_i\xi,
 \qquad \xi\in\{1,t,z,w\},\quad a,b\in\mathbb Z.
\end{equation}
The data also multiply PSD-based or ratio terms by explicitly specified products of $t,z,w$. The exact term names determine which products and integers are used; all are evaluated with $w=t+z-18$. Every factor is nonnegative under the hypotheses: for the ratio factor use Lemma~\ref{lem:ratio}, and for the squared factors no sign assumption on $h$ or $q$ is necessary.

The ancillary checker reconstructs the identity exactly, including the substitutions~\eqref{eq:mixed}, and confirms equality of every rational coefficient. The fixed Lean module \texttt{APPT.Uniform} proves~\eqref{eq:uniform-nonneg} using the same exact data; its dependencies supply evaluation-preserving integer coefficient operations and sign proofs. With $T=1$, divide by the positive number $8(9+t+z)$ to obtain~\eqref{eq:uniform-bound}.
\end{proof}

\begin{proposition}[All large dimensions]\label{prop:large}
For every integer $n\geq9$, every APPT density matrix on $\C^3\otimes\C^n$ has purity at most $3/(8n)$.
\end{proposition}
\begin{proof}
Let $D=3n$, choose $y$ as in~\eqref{eq:outer}, and let $x$ contain the $M=D-9$ middle entries $\lambda_4,\ldots,\lambda_{D-6}$. Then $y_4\leq x_i\leq y_3$ and $M\geq18$. By Lemma~\ref{lem:compression}, there are $t,z\geq0$ with $t+z=M$ for which
\[
 T(y,t,z)=\sum_{i=1}^D\lambda_i=1,\qquad
 \sum_{i=1}^D\lambda_i^2\leq S(y,t,z).
\]
Lemma~\ref{lem:necessity} supplies $A(y),B(y)\succeq0$. The uniform certificate therefore gives
\[
 \Tr(\rho^2)=\sum_{i=1}^D\lambda_i^2
 \leq\frac9{8(9+t+z)}=\frac9{8D}=\frac3{8n}.
\]
The argument treats every integer $n\geq9$ through one real-parameter identity; it is not an extrapolation from finitely many dimensions.
\end{proof}

\section{Actual APPT attaining states}\label{sec:attainment}
We provide physical witnesses without relying on the sufficiency direction of a spectral APPT criterion.

\begin{lemma}[A rank-one perturbation]\label{lem:rankone}
For every unit vector $v\in\C^3\otimes\C^n$, the matrix $I+2vv^*$ is APPT before trace normalization.
\end{lemma}
\begin{proof}
A unit vector has Schmidt coefficients $s_1,s_2,s_3\geq0$ with $\sum_i s_i^2=1$, allowing zero coefficients. The nonzero eigenvalues of $(vv^*)^\Gamma$ are $s_i^2$ and $\pm s_i s_j$ for $i<j$. Since $2s_i s_j\leq s_i^2+s_j^2\leq1$, we have $I+2(vv^*)^\Gamma\succeq0$. For every global unitary $U$, the vector $Uv$ is again a unit vector, so the same argument applies to $U(I+2vv^*)U^*=I+2(Uv)(Uv)^*$. The formalization proves the same rank-one bound by an explicit skew-Gram identity and a rectangular-contraction argument instead of assuming a Schmidt-decomposition interface.
\end{proof}

\begin{lemma}[A positive contraction]\label{lem:contraction}
If $0\preceq P\preceq I$ on $\C^3\otimes\C^n$, then $P^\Gamma+I\succeq0$, and $I+P$ is APPT before normalization.
\end{lemma}
\begin{proof}
Let $E_{ij}$ be the matrix units on $\C^3$, let $D_i=E_{ii}$, and for $i<j$ put $S_{ij}=E_{ij}+E_{ji}$ and $K_{ij}=E_{ji}-E_{ij}$. Tensor these matrices with $I_n$ when they act on the full space. Direct block multiplication gives the exact identity
\begin{equation}\label{eq:contraction-identity}
 P^\Gamma+I=\sum_{i=1}^3D_iPD_i^*
 +\frac12\sum_{i<j}\bigl(S_{ij}PS_{ij}^*+K_{ij}(I-P)K_{ij}^*\bigr).
\end{equation}
Every summand is PSD. For a global unitary $U$, the conjugate $Q=UPU^*$ is another positive contraction, so $I+Q^\Gamma\succeq0$. This is exactly the APPT condition for $I+P$.
\end{proof}

\begin{proof}[Proof of Theorem~\ref{thm:main}]
The upper bounds are Corollary~\ref{cor:finite} and Proposition~\ref{prop:large}. For the first witness, $R=vv^*$ satisfies $R^2=R$ and $\Tr R=1$. Hence $I+2R\succeq0$, $\Tr(I+2R)=D+2$, and
\[
 \Tr\left(\left(\frac{I+2R}{D+2}\right)^2\right)
 =\frac{\Tr(I+8R)}{(D+2)^2}=\frac{D+8}{(D+2)^2}.
\]
Lemma~\ref{lem:rankone} gives APPT. For the second witness choose a rank-$n$ orthogonal projection $P$, for example the coordinate projection onto $e_1\otimes\C^n$. Then $P^2=P$, $\Tr P=n$, and Lemma~\ref{lem:contraction} gives APPT. Its trace normalization is $D+n=4n$, and
\[
 \Tr\left(\left(\frac{I+P}{4n}\right)^2\right)
 =\frac{\Tr(I+3P)}{16n^2}=\frac{6n}{16n^2}=\frac3{8n}.
\]
Positive scalar normalization preserves APPT. Thus each upper bound is attained by a genuine density matrix.
\end{proof}

\begin{corollary}[Circumradius]\label{cor:radius}
The squared radius of the smallest Hilbert--Schmidt ball centered at $I/(3n)$ containing the APPT density matrices is
\[
 R_n^2=\begin{cases}
 \displaystyle\frac{4(3n-1)}{3n(3n+2)^2},&3\leq n\leq8,\\[5pt]
 \displaystyle\frac1{24n},&n\geq9.
 \end{cases}
\]
\end{corollary}
\begin{proof}
For a Hermitian trace-one matrix, $\|\rho-I/D\|_\hs^2=\Tr(\rho^2)-1/D$. Substitute Theorem~\ref{thm:main}. This is an algebraic corollary of the formalized purity result; the release does not export a separately named Hilbert--Schmidt circumradius theorem.
\end{proof}


\section{A higher-dimensional counterexample}
\label{sec:counterexample}
The general candidate~\eqref{eq:general-candidate} is the purity maximum of the inscribed spectral polytope in~\cite{AKW}; Conjecture~6.7 there asserts that it is also the maximum over all APPT spectra. We now give an actual APPT state outside that polytope whose purity is larger.

\begin{theorem}[Explicit failure of the general formula]\label{thm:counterexample}
On $\C^{10}\otimes\C^{38}$, the diagonal state
\begin{equation}\label{eq:explicit-counterexample}
 \rho_{10,38}=\frac{\operatorname{diag}(41,\underbrace{17,\ldots,17}_{146},
 \underbrace{15,\ldots,15}_{233})}{6018}
\end{equation}
is APPT and satisfies
\[
 \Tr(\rho_{10,38}^2)=\frac{2675}{1006009}>M_{10,38}.
\]
More precisely, the two candidate values are $97/36481$ and $5/1881$, and the exact gaps are respectively
\[
 \frac{3802}{36700214329},\qquad \frac{1630}{1892302929}.
\]
Consequently Conjecture~6.7 of~\cite{AKW} is false in its stated arbitrary-dimensional form.
\end{theorem}

\begin{proof}
We prove a slightly more general all-unitary lemma. Let $m\geq4$, $n\geq m$, $D=mn$, let $P$ be an orthogonal projection on $\C^m\otimes\C^n$, and let $v$ be a unit vector in the range of $P$. Set
\[
 A=(2m-5)I+2P+4(m-4)\lvert v\rangle\langle v\rvert.
\]
For a unit vector $\psi$ with Schmidt coefficients $s_1\geq\cdots\geq s_m\geq0$, put $W=(\lvert\psi\rangle\langle\psi\rvert)^\Gamma$. The negative eigenvalues of $W$ are $-s_i s_j$ for $i<j$, and its least eigenvalue is $-s_1s_2$. Since $0\leq UPU^*\leq I$, spectral evaluation gives
\[
 \Tr(UPU^*W)\geq-\sum_{i<j}s_i s_j,
 \qquad
 \langle Uv,WUv\rangle\geq-s_1s_2.
\]
Thus, for every global unitary $U$,
\begin{align*}
 \langle\psi,(UAU^*)^\Gamma\psi\rangle
 &\geq (2m-5)\sum_i s_i^2-2\sum_{i<j}s_i s_j
       -4(m-4)s_1s_2\\
 &= (2m-6)(s_1-s_2)^2
 +2\sum_{3\leq i<j\leq m}(s_i-s_j)^2
 +(s_1+s_2-\sum_{i=3}^m s_i)^2\geq0.
\end{align*}
The identity is obtained by expanding both sides, so $A$ is APPT. If $\operatorname{rank}P=k$, the eigenvalues of $A$ are
\[
 6m-19\ (1\text{ time}),\qquad 2m-3\ (k-1\text{ times}),
 \qquad 2m-5\ (D-k\text{ times}).
\]
For $m=10,n=38,k=147$ this is exactly~\eqref{eq:explicit-counterexample}. Its trace-square is $96300/6018^2$, which reduces to $2675/1006009$. Substitution of $m=10,n=38,t=171$ into~\eqref{eq:general-candidate} gives the two displayed candidates and the stated positive rational differences. No numerical optimization or sampled unitary is used.
\end{proof}

The counterexample is not confined to a singular boundary. If $\varepsilon=1/1000$ and
\[
 \rho_\varepsilon=(1-\varepsilon)\rho_{10,38}+\varepsilon I/380,
\]
then every global-unitary partial transpose is bounded below by $I/380000$, while its purity still exceeds both candidates. The same construction gives an explicit infinite family for every $m\geq11$ with $n=4m$; the exact factorization and multiplicities are included in the supplementary proof file.

\section{The corrected unrestricted asymptotic law}
\label{sec:asymptotic}
The finite counterexample leads to a sharp asymptotic replacement for the false formula. The result concerns the regime in which the smaller local dimension grows; it does not assert an exact finite-dimensional maximum.

\begin{theorem}[Uniform APPT excess-purity asymptotic]\label{thm:asymptotic}
Let $P_{\max}(m,n)$ be the maximum purity over APPT density matrices on $\C^m\otimes\C^n$. Then, as $m\to\infty$, uniformly over integers $n\geq m$,
\begin{equation}\label{eq:asymptotic-law}
 \sup_{n\geq m}\left|
 \frac{D^2\left(P_{\max}(m,n)-1/D\right)}{\max\{8,4+n/m\}}-1\right|\longrightarrow0,
 \qquad D=mn.
\end{equation}
In particular, for every fixed $\gamma\geq1$ and every sequence $n/m\to\gamma$,
\[
 D^2\left(P_{\max}(m,n)-\frac1D\right)\longrightarrow\max\{8,4+\gamma\},
\]
and $P_{\max}(m,m)=m^{-2}+8m^{-4}+o(m^{-4})$.
\end{theorem}

\begin{proof}
We give the two estimates that are uniform in the spectrum and in $n$. Order the eigenvalues of an APPT state as $\lambda_1\geq\cdots\geq\lambda_D$, put $R=m(m-1)/2$, $S=m(m+1)/2$, $j=D-S+1$, $b=\lambda_j$, and $\delta_i=\lambda_i-\lambda_{D+1-i}$. For every assignment of the nonnegative $\delta_i$ to the edges of $K_m$, a Schmidt-vector partial-transpose test gives
\begin{equation}\label{eq:edge-test}
 \lambda_{\max}(A_\delta)\leq2b,
\end{equation}
where $A_\delta$ is the resulting weighted adjacency matrix. This uses the top $R$ eigenvalues on antisymmetric pair vectors, the bottom $R$ on symmetric pair vectors and the next $m$ eigenvalues on diagonal vectors; those slots are distinct because $m^2\leq D$.

Place the largest $m-1$ gaps on a star. Equation~\eqref{eq:edge-test} gives $\sum_{i=1}^{m-1}\delta_i^2\leq4b^2$. In a second, row-by-row placement, subtract the smallest gap $w=\delta_R$ from every edge and denote the residuals by $\eta_i=\delta_i-w$. The elementary row-degree inequality
\[
 \sum_{i=m}^R\eta_i^2\leq\|A_\eta m^{-1/2}{\bf1}\|^2
\]
together with~\eqref{eq:edge-test} yields
\[
 V:=\sum_i\lambda_i^2-1/D\leq C(m,n)b^2,
 \qquad C(m,n)=\max\left\{8,4+\frac{D}{(m-1)^2}\right\}.
\]
The derivation uses the midpoint of $\lambda_R$ and $\lambda_{D-R+1}$; all remaining eigenvalues lie within $w/2$ of that midpoint. Since $jb\leq1$ and the first $j$ eigenvalues contribute at least $j(b-1/D)^2$ to $V$ whenever $b>1/D$, this gives
\[
 D^2V\leq\frac{C(m,n)}{(1-\sqrt{C(m,n)/j})^2}.
\]
As $m\to\infty$, the denominator tends uniformly to one and $C(m,n)\leq[m/(m-1)]^2\max\{8,4+n/m\}$.

For the matching lower bound, two explicit APPT constructions are used. A one-spike/plateau state gives, for every finite $m\geq4$ and $n\geq m$, a centered coefficient tending uniformly to $4+n/m$. A finite lacunary hierarchy of nested spectral plateaus gives a coefficient arbitrarily close to $8$, with a fixed positive slack in the all-unitary graph bound before the dimension limit. The hierarchy is first held fixed, then $m$ tends to infinity, and only afterward are the slack and number of levels optimized. Thus for every $\varepsilon>0$, all sufficiently large $m$ admit APPT states with
\[
 D^2(P_{\max}-1/D)\geq(1-\varepsilon)\max\{8,4+n/m\}
\]
uniformly in $n\geq m$. Combining the upper and lower estimates proves~\eqref{eq:asymptotic-law}. The complete graph-limit and forest argument for the lacunary construction is written in Appendices~\ref{app:sharp-upper}--\ref{app:rectangular}, with exact rational regression checks; no finite numerical test is used as a premise.
\end{proof}

\begin{remark}[Status of the higher-dimensional argument]
Theorem~\ref{thm:asymptotic} is a written analytic proof supported by exact rational and integer checkers. It is not part of the qutrit Lean release. It does not determine exact finite-dimensional maxima, classify maximizing spectra, or prove equality between APPT and absolutely separable states.
\end{remark}

\begin{theorem}[All-regime equivalent]\label{thm:allregime}
Let $P_{\max}(m,n)$ be as above and define
\[
 A(m,n)=\frac{\max\{8,4+D/(m^2-1)\}}{D^2},\qquad D=mn.
\]
Then
\begin{equation}\label{eq:allregime}
 \lim_{N\to\infty}\ \sup_{2\leq m\leq n,\,mn\geq N}
 \left|\frac{P_{\max}(m,n)-1/(mn)}{A(m,n)}-1\right|=0.
\end{equation}
For each fixed $m\geq2$, this implies
\[
 D\,P_{\max}(m,n)\longrightarrow\frac{m^2}{m^2-1},
 \qquad
 P_{\max}(m,n)-\frac1D\sim\frac1{(m^2-1)D}
 \quad(n\to\infty).
\]
\end{theorem}

\begin{proof}
Theorem~\ref{thm:asymptotic} proves the uniform statement when the smaller dimension tends to infinity. It remains to treat a bounded subsequence of $m$. Fix $m\geq2$ and use the notation $R=m(m-1)/2$, $S=m(m+1)/2$, $j=D-S+1$, and $b=\lambda_j$ for an ordered APPT spectrum. The maximally entangled Schmidt-rank-$m$ test, with the largest $R$ eigenvalues placed on its negative witness space, gives
\[
 \sum_{i=1}^R\lambda_i\leq\sum_{i=D-S+1}^{D}\lambda_i,
\]
so $\lambda_R\leq[(m+1)/(m-1)]b$. A single-edge test gives $\lambda_1-\lambda_D\leq2b$, hence $\lambda_1\leq3b$. The $N_0=D-m^2+2$ eigenvalues from indices $R$ through $D-S+1$ therefore lie in $[b,rb]$, where $r=(m+1)/(m-1)$. For numbers in this interval,
\[
 x^2\leq(r+1)bx-rb^2.
\]
After summing and optimizing the resulting quadratic in $z=N_0b/T$, one obtains
\[
 \sum_i\lambda_i^2\leq\frac{m^2/(m^2-1)}{D-m^2+2}
 +\frac{9(m^2-2)}{(D-S+1)^2}.
\]
This gives the fixed-$m$ upper limit after multiplication by $D$. For the matching lower bound, let $c=2/(m-1)$, choose a rank-$k$ projection with $k=\lceil D/(c+2)\rceil$, and normalize $I+cP$. The sum of the negative Schmidt eigenvalues is at most $(m-1)/2$, so this state is APPT for every global unitary. Its exact purity calculation gives
\[
 D\,\Tr(\rho^2)\longrightarrow\frac{(c+2)^2}{4(c+1)}=\frac{m^2}{m^2-1}.
\]
Thus the fixed-$m$ limit follows. Finally, if~\eqref{eq:allregime} failed, a sequence with $D\to\infty$ and a fixed error would have either a bounded subsequence of $m$ (handled by the fixed-$m$ argument) or $m\to\infty$ (handled by Theorem~\ref{thm:asymptotic}). In the latter case $D/(m^2-1)$ and $n/m$ differ by a factor tending to one. This sequential criterion proves~\eqref{eq:allregime}. The detailed fixed-$m$ derivation and exact deficit identity are recorded in Appendix~\ref{app:uniform-lower}.
\end{proof}

\section{An exact rectangular regime}
\label{sec:rectangular}
The asymptotic theorem has a finite-dimensional companion in sufficiently rectangular systems. The cutoff below is a proved sufficient threshold; it is not asserted to be the smallest threshold for the actual APPT problem.

\begin{theorem}[Exact maximum and all maximizers in a rectangular regime]\label{thm:rectangular}
Let $3\leq m\leq n$ and $D=mn$. If
\[
 n\geq n_0(m):=m^3-m-2,
 \qquad t=\left\lceil\frac{(m-1)n}{2}\right\rceil,
\]
then
\begin{equation}\label{eq:rectangular-max}
 P_{\max}(m,n)=\frac{D(m-1)^2+4mt}{[D(m-1)+2t]^2}.
\end{equation}
Every maximizing density matrix has spectrum
\begin{equation}\label{eq:rectangular-spectrum}
 \frac{(\underbrace{m+1,\ldots,m+1}_{t},
 \underbrace{m-1,\ldots,m-1}_{D-t})}{D(m-1)+2t},
\end{equation}
so all maximizers form the global unitary orbit of
\[
 \frac{(m-1)I+2P}{D(m-1)+2t},
 \qquad \operatorname{rank}P=t.
\]
\end{theorem}

\begin{proof}
We use one necessary physical test and optimize its outer polytope exactly. Put $R=m(m-1)/2$ and $S=m(m+1)/2$. Applying the maximally entangled Schmidt-rank-$m$ test after a global unitary gives, for every decreasing APPT spectrum,
\begin{equation}\label{eq:outer-single-test}
 \sum_{i=1}^{R}\lambda_i\leq\sum_{i=D-S+1}^{D}\lambda_i.
\end{equation}
Let $\mathcal O_{m,n}$ be the ordered probability simplex cut out by~\eqref{eq:outer-single-test}. Its gap coordinates give the decomposition into prefix-uniform vectors $u_i=(1/i,\ldots,1/i,0,\ldots,0)$. The cut has one sign change at $L=D-m$, so every vertex of $\mathcal O_{m,n}$ is either a surviving $u_j$ with $j\geq L$ or a boundary point on an edge joining $u_i$ with $i<L<j$. This follows directly from the two affine constraints: any point with three positive gap coordinates admits a nontrivial perturbation preserving them and hence is not a vertex.

Writing $j=L+s$ with $0\leq s\leq m$ and using the exact boundary interpolation, the edge vertices have the explicit spectrum in Appendix~\ref{app:rectangular}. The early, middle, and late index ranges reduce to
\[
 \max_{\lambda\in\mathcal O_{m,n}}\sum_i\lambda_i^2
 =\max\left\{Q_{\rm sp},Q_t\right\},
\]
where
\[
 Q_{\rm sp}=\frac{D+m(m+2)}{(D+m)^2},\qquad
 Q_t=\frac{D(m-1)^2+4mt}{[D(m-1)+2t]^2}.
\]
In the middle range, differentiating the full-support endpoint in the integer index shows its unique maximum at $t=\lceil(m-1)n/2\rceil$, including the odd-parity case. Exact subtraction shows $Q_t>Q_{\rm sp}$ whenever $n\geq m^3-m-2$. The complete vertex calculation, including the three index ranges and all rounding cases, is proved in Appendix~\ref{app:rectangular}; the checker additionally verifies its algebra in exact rational arithmetic.

It remains to show that the winning outer vertex is genuinely APPT. For a unit Schmidt vector with coefficients $s_i$, the sum of the absolute values of the negative eigenvalues of its partial transpose is $\sum_{i<j}s_i s_j\leq(m-1)/2$. Thus, for every conjugated rank-$t$ projection $P'$, $\operatorname{Tr}(P'W)\geq-(m-1)/2$, where $W$ is the partial transpose of the Schmidt projector. Hence $((m-1)I+2P')^\Gamma\succeq0$ for every global unitary, proving attainment of $Q_t$.

Finally, the squared Euclidean norm is strictly convex on the ordered simplex. Since the winning vertex is unique, no nontrivial convex combination of other vertices can attain the same value. Every APPT spectrum lies in $\mathcal O_{m,n}$, so its equality spectrum is exactly~\eqref{eq:rectangular-spectrum}; Hermitian matrices with that spectrum are precisely its global unitary orbit. This proves both the maximum and the equality classification.
\end{proof}

\begin{remark}[Scope of the rectangular theorem]
The cutoff $n_0(m)$ is a sufficient threshold for the actual APPT maximum and is the exact transition for the single-test outer-polytope argument used here. It is not claimed to be the smallest dimension at which the two-level spectrum is globally optimal. The explicit $10\times38$ counterexample lies below its cutoff and therefore does not conflict with this theorem. The proof is analytic and its finite vertex calculations are supported by exact rational checks, not Lean formalization.
\end{remark}

\section{Qutrit certificate data and complete formal verification}\label{sec:verification}
\subsection{Data specification and independent checking}
The ancillary directory contains all seven JSON files listed in Appendix~\ref{app:data}, copied byte for byte from the fixed formalization. Each contains a \texttt{terms} list. Its entries have a \texttt{coefficient} field holding a rational string and a \texttt{name} field specifying the polynomial. The checker reconstructs the tail sums, matrices~\eqref{eq:A}--\eqref{eq:B}, their determinants, principal minors and quadratic forms. It then expands~\eqref{eq:finite-cert} or~\eqref{eq:uniform-cert} and verifies coefficient equality. The reconstruction code uses only exact integers and rational numbers; no optimizer is imported. The included input marker \texttt{exact\_certificate} is checked but is not accepted as proof: coefficient equality is recomputed.

The JSON term language uses zero-based indices in explicit exponent lists and \texttt{gprod} tuples, and one-based subscripts in names such as \texttt{Aminor02*g23}; the latter means the principal minor on matrix indices $0,2$ multiplied by the twenty-third gap. The supplied reconstruction code is the unambiguous machine-readable specification of these conventions. In the uniform file, \texttt{param} entries use $0,1,2$ to denote $t,z,t+z-18$. All term signs are justified as in the two certificate lemmas; there is no search over or sampling of real parameter values in the verification.

From the preprint source root, the independent rational check is
\begin{verbatim}
python3 ancillary/check_certificates.py
\end{verbatim}
The complete formalization is available separately in the immutable release \cite{ProofRelease}. Its source manifest fixes both the coefficient data and the Lean proofs. This separation keeps the preprint source archive small while preserving a precise reference to every formal dependency.

\subsection{Formal statement and dependency map}
The final Lean theorem in namespace \texttt{APPT.Quantum} has the form
\begin{verbatim}
theorem appt_purity_maximum_formula (n : Nat) (hn : 3 <= n) :
  IsGreatest (attainablePurities n)
    (if n <= 8 then (3*(n : Real)+8)/(3*(n : Real)+2)^2
     else 3/(8*(n : Real)))
\end{verbatim}
Here \texttt{attainablePurities} quantifies over actual matrices indexed by \texttt{Fin 3 $\times$ Fin n}; \texttt{IsDensity} means PSD and complex trace one; \texttt{AbsolutelyPPT} quantifies over the actual complex unitary group; and \texttt{purity A} is the real part of the actual trace of \texttt{A*A}. For a density matrix the spectral-moment theorem identifies this real quantity with the usual real purity.

\begin{center}\small
\begin{tabular}{>{\raggedright\arraybackslash}p{0.38\textwidth}>{\raggedright\arraybackslash}p{0.53\textwidth}}
\toprule
Proof obligation&Principal Lean modules or declarations\\
\midrule
Global-unitary corner necessity&\texttt{Quantum.CornerNecessity}, \texttt{Quantum.SpectralNecessity}\\
Ordered eigenvalues and moments&\texttt{Quantum.SpectralData}, \texttt{Quantum.SpectralMoments}\\
Six finite certificate bounds&\texttt{Finite9Bound} through \texttt{Finite24Bound}, and \texttt{SpectrumBound*}\\
Uniform and middle-spectrum bound&\texttt{Uniform}, \texttt{Compression}, \texttt{UniformBridge}, \texttt{OrderedSpectrum}\\
Physical APPT witnesses&\texttt{Quantum.ShortAttainment}, \texttt{Quantum.LongAttainment}, \texttt{Quantum.Attainment}\\
All-$n$ maximum&\texttt{Quantum.Maximum}: \texttt{appt\_purity\_maximum\_formula}\\
\bottomrule
\end{tabular}
\end{center}

\subsection{Build, replay, and rejection controls}
The proof uses Lean 4.34.1 and Mathlib commit
\begin{center}\small\texttt{d13f23b723b8a846827a245b89c10fc7d3f11612}.
\end{center}
The published proof commit is \proofcommit. The final successful verification run is \texttt{37968137659}, at verified source commit
\begin{center}\small\texttt{f372047c63ba501b13301b4b3ba089c261e3722e}.
\end{center}
All 942 pinned proof, tool and data files agree between the verified and published sources. The mathematical theorem files are unchanged from the first successful full proof replay; later corrections concerned auxiliary tests and verification orchestration.

All 879 local modules and the default Lake target passed. The final audit replayed the transitive closure of 58 roots, including the exact maximum theorem: 155,787 declarations were checked in an initially empty Lean kernel at trust level zero. The only axioms in the audited closure are \texttt{propext}, \texttt{Classical.choice}, and \texttt{Quot.sound}. The checker rejects unsafe or partial mathematical dependencies and unapproved axioms, and compares each replayed root's type and universe parameters with the original. No \texttt{sorry}, \texttt{admit}, custom axiom or untrusted \texttt{native\_decide} is used in the mathematical closure.

The paired failure controls are part of the record. They alter a coefficient, a PSD-corner example, or an endpoint value, and require the expected rejection rather than an arbitrary process failure. A further test replaces only the proof of the same final maximum theorem by \texttt{True.intro}, preserving its name, type and universe parameters. The changed declaration's 8,581 actual dependencies are positively checked first; a fresh trust-zero kernel then rejects the forged theorem for a declaration type mismatch. The large original proof has already passed the full positive replay. Dependencies belonging only to the discarded original proof are not unnecessarily repeated in the negative test.

The final workflow reused independently compiled components bound to exact source, object and log hashes. It does \emph{not} claim a fresh build of all 879 modules on a single runner. This does not change the independent empty-kernel replay of the full mathematical closure. Each recorded local-module build had a 300-second limit; the largest recorded time was 239.703 seconds. Small coefficient lemmas, literal intermediate coefficient lists and explicit integer constructors avoid an unbounded monolithic elaboration.

The fresh reproduction command, from the formalization's package root, is
\begin{verbatim}
lake exe cache get
python3 scripts/source_manifest.py --check
python3 reproduce.py --fresh --jobs 2
\end{verbatim}
The fresh option removes the project's own build directory, not the external Mathlib caches. Literal outputs, source hashes, component provenance, dependency inventories and the earlier failed audit attempts are retained. A failed historical workflow is never relabeled as a successful one.

\section{An analytic corollary for absolute separability}\label{sec:asep}
A state is \emph{absolutely separable} if every global-unitary conjugate is separable across the fixed bipartition. Since separable states are PPT, the absolutely separable set is contained in the APPT set. The following consequence concerns its extremal purity, not equality of the two sets.

\begin{corollary}[Absolute-separability purity]\label{cor:asep}
For every $n\geq3$, the maximum purity of absolutely separable states on $\C^3\otimes\C^n$ equals $p_n$ in~\eqref{eq:maximum}.
\end{corollary}
\begin{proof}
The APPT theorem supplies the upper bound. For the large branch, use the sufficient condition of Kondra, Barrios Hita, Neumann, Kampermann and Bru{\ss} \cite[Supplemental Material, Lemma 2]{Kondra}: in local dimensions $m\leq n$, a full-rank state with
\[
 \lambda_{\max}/\lambda_{\min}\leq(m+1)/(m-1)
\]
is separable. As the criterion is spectral, it also certifies absolute separability. The state $(I+P)/(4n)$ has eigenvalue ratio $2$, exactly this threshold for $m=3$.

For the small branch, we give a direct decomposition valid for every unit vector $v$. Use local unitaries to write $v=\sum_{i=1}^r s_i|ii\rangle$, with $s_i\geq0$ and $\sum_i s_i^2=1$. Let the independent phases $\theta_i$ be uniform on $\{0,2\pi/3,4\pi/3\}$ and set
\[
 x_\theta=\sum_i\sqrt{s_i}e^{\mathrm{i}\theta_i}|i\rangle,
 \qquad z_\theta=\sum_i\sqrt{s_i}e^{-\mathrm{i}\theta_i}|i\rangle.
\]
A finite phase average of product projectors gives
\[
 \Omega=\mathbb E_\theta\bigl[|x_\theta\otimes z_\theta\rangle
        \langle x_\theta\otimes z_\theta|\bigr]
 =vv^*+\sum_{i\ne j}s_i s_j|ij\rangle\langle ij|.
\]
Indeed, the only nonzero averaged phase factors correspond to identical ordered pairs or to two diagonal pairs. Consequently
\[
 I+2vv^*=2\Omega+\sum_i|ii\rangle\langle ii|
       +\sum_{i\ne j}(1-2s_i s_j)|ij\rangle\langle ij|+I_{\mathrm{out}},
\]
where the displayed sums range over the Schmidt-support product subspace and $I_{\mathrm{out}}$ is the identity on its complementary product-basis coordinates. Every term is separable and every coefficient is nonnegative because $2s_i s_j\leq s_i^2+s_j^2\leq1$. Thus $(I+2vv^*)/(3n+2)$ is separable for every unit $v$. Any global-unitary conjugate has the same form with a new unit vector, proving absolute separability and the required attainment.
\end{proof}

\begin{remark}[Formalization boundary]
Corollary~\ref{cor:asep} is an analytic consequence using the explicitly cited separability criterion and the decomposition above. It is \emph{not} included in the immutable Lean proof of the APPT theorem. No claim is made that the separability cone, the cited criterion, or this corollary has been formalized in that release. In particular, equal maximal purities do not imply equal convex sets.
\end{remark}

\section{Scope and further questions}
The main theorem closes the maximal-purity problem for $3\otimes n$ with $n\geq3$. The $3\otimes2$ system is, after swapping factors, a qubit--qutrit problem and is not covered by the present formal theorem. The immutable qutrit proof does not classify its maximizers. The new analytic rectangular theorem separately classifies all maximizers in its stated range, and the higher-dimensional counterexample refutes the general formula. Corollary~\ref{cor:asep} gives the analogous absolute-separability maximum analytically; the immutable Lean release covers only the APPT theorem.

The arbitrary-dimensional inscribed-polytope purity conjecture of \cite{AKW} is false by Theorem~\ref{thm:counterexample}. Theorems~\ref{thm:asymptotic} and~\ref{thm:allregime} give the corrected first excess-purity term whenever the total dimension grows, and Theorem~\ref{thm:rectangular} gives the exact finite maximum and all maximizers beyond a proved rectangular threshold. The remaining finite-dimensional regimes and a separate, more structural question---whether absolute PPT and absolute separability coincide beyond qubit--qudit systems---remain open. The additional witness analysis in Section~\ref{sec:asep} establishes a common qutrit maximum, but proves no reverse inclusion between the two sets. Equal purity maxima must not be mistaken for equality of the convex sets.

\paragraph{Author, funding and AI assistance.}
This research was conducted independently and received no external funding. The institutional affiliation identifies the author's place of study and does not assert institutional sponsorship or endorsement. AI assistance was used in mathematical exploration, certificate and program development, Lean proof development, verification orchestration and manuscript drafting. Formal acceptance refers to the stated Lean declarations and audited dependency closure; it is not a claim of external peer review or independent human certification of every narrative statement. No separate claim of historical priority is made beyond the dated public artifacts and the explicit comparison with prior work above.

\paragraph{Data and code availability.}
The exact qutrit data and independent checker accompany this preprint. The complete qutrit proof, reproduction scripts and verification evidence are frozen in \cite{ProofRelease}; the earlier immutable interim release is left unchanged. The higher-dimensional analytic proof and its standard-library exact checkers are included in the supplementary directory, with the source files and hashes recorded in the repository. This manuscript is a preprint, not a journal-accepted article.

\appendix
\section{The unrestricted upper estimate}\label{app:sharp-upper}
This appendix supplies the details used in Theorem~\ref{thm:asymptotic}.
Throughout, $m\geq3$, $n\geq m$, $D=mn$, and the eigenvalues
$\lambda_1\geq\cdots\geq\lambda_D\geq0$ sum to one. Set
\[
 R=\frac{m(m-1)}2,\quad S=\frac{m(m+1)}2,\quad j=D-S+1,\quad
 b=\lambda_j,\quad \delta_i=\lambda_i-\lambda_{D+1-i}\quad(1\leq i\leq R).
\]
The gaps $\delta_i$ are nonnegative and decreasing.

\subsection{Every edge rearrangement is a physical test}
For $x_i\geq0$ with $\sum_i x_i^2=1$, the partial transpose $W$ of the
projector onto $\sum_i x_i|ii\rangle$ has eigenvalues $x_i^2$ on diagonal
vectors, $+x_ix_k$ and $-x_ix_k$ on symmetric and antisymmetric pair
vectors, and zero on the remaining $D-m^2$ dimensions. For any assignment
of the $R$ gaps to the edges of $K_m$, assign $\lambda_i$ to the negative
pair vector of its edge and $\lambda_{D+1-i}$ to the positive pair vector.
Assign the next $m$ smallest eigenvalues to the diagonal vectors. These
are all at most $b$, and all slots are distinct since $R+S=m^2\leq D$.
A global unitary realizes the assignment. If $A_\delta$ is the
zero-diagonal symmetric matrix of edge weights, the APPT condition gives
\[
 0\leq\sum_i\beta_i x_i^2-\tfrac12x^TA_\delta x
 \leq b-\tfrac12x^TA_\delta x.
\]
A symmetric nonnegative matrix has a nonnegative maximizing Rayleigh
vector: replacing each coordinate by its absolute value can only
increase the quadratic form. Moreover
$|z^TA_\delta z|\leq |z|^TA_\delta|z|$. Thus
\begin{equation}\label{eq:all-edge-norm}
 \|A_\delta\|_{\rm op}\leq2b
\end{equation}
for every rearrangement. The all-ones vector also gives
$\sum_i\delta_i\leq mb$.

\subsection{A triangular degree estimate}
For any decreasing list $\eta_1,\ldots,\eta_R\geq0$, place the entries in
the upper triangle row by row and let $A_\eta$ be the symmetric
adjacency matrix. With $e=m^{-1/2}{\bf1}$,
\begin{equation}\label{eq:triangular}
 \sum_{i=m}^{R}\eta_i^2\leq\|A_\eta e\|^2.
\end{equation}
Indeed, let $a_t$ be the smallest weight in upper-triangular row $t$,
for $1\leq t\leq m-2$. Every incoming edge at vertex $t$ and every
outgoing edge of that row has weight at least $a_t$, so the degree
at that vertex is at least $(m-1)a_t$. Row $t+1$ has $m-t-1$ weights,
all at most $a_t$. Since $(m-1)^2/m\geq m-2\geq m-t-1$, that row's
square sum is bounded by the squared degree of vertex $t$ divided
by $m$. Summing over $t$ proves~\eqref{eq:triangular}.

\subsection{Variance, background and normalization}
Put $w=\delta_R$, $\eta_i=\delta_i-w$, and $A=\sum_i\eta_i$.
Place the largest $m-1$ original gaps on a star and fill all other
edges with the remaining weights. The star's norm is the square root
of their square sum, and its nonnegative Rayleigh vector is also a
test for the full graph. Hence
\[
 \sum_{i=1}^{m-1}\eta_i^2\leq4b^2.
\]
For a different rearrangement, use the triangular placement of
$\eta$ and add $w$ to every edge. Expanding its squared norm at $e$
in~\eqref{eq:all-edge-norm} gives
\[
 \|A_\eta e\|^2+\frac{4(m-1)}m wA+(m-1)^2w^2\leq4b^2,
\]
because the sum of the residual degrees is $2A$. It follows that
\begin{equation}\label{eq:residual-square}
 \sum_i\eta_i^2\leq8b^2-\frac{4(m-1)}m wA-(m-1)^2w^2.
\end{equation}
Let $z=(\lambda_R+\lambda_{D-R+1})/2$. In the $i$th top/bottom pair,
$u=\lambda_i-z$ and $v=z-\lambda_{D+1-i}$ are at least $w/2$ and sum
to $\delta_i$. Expanding $(u-w/2)(v-w/2)\geq0$ yields
$u^2+v^2\leq\delta_i^2-w\delta_i+w^2/2$.
The other $D-2R$ eigenvalues lie within $w/2$ of $z$.
Centering at the mean can only reduce the squared distance, so
\[
 V=\sum_i\lambda_i^2-\frac1D
 \leq\sum_i\eta_i^2+wA+\frac D4w^2.
\]
Insert~\eqref{eq:residual-square}. The coefficient of $wA$ is
$1-4(m-1)/m\leq0$, and the all-ones test gives
$(m-1)w\leq2b$. Maximizing the remaining quadratic over that interval
proves the finite inequality
\begin{equation}\label{eq:finite-sharp-upper}
 V\leq Cb^2,\qquad C=\max\{8,4+D/(m-1)^2\}.
\end{equation}
There is no assumption on the number of spectral levels.
Also $jb\leq1$ and $b\leq1/D+\sqrt{V/j}$: when $b>1/D$ the first
$j$ eigenvalues contribute at least $j(b-1/D)^2$ to $V$.
If $C<j$, solving~\eqref{eq:finite-sharp-upper} for $\sqrt V$ gives
\[
 D^2V\leq\frac{C}{(1-\sqrt{C/j})^2}.
\]
For all $n\geq m\geq3$, $j\geq D/3$ and
$C/j\leq24/m^2+3/(m-1)^2$. Writing $K=\max\{8,4+n/m\}$,
we have $C\leq[m/(m-1)]^2K$. These estimates prove the uniform
upper bound in Theorem~\ref{thm:asymptotic}.

\section{Lacunary plateaus and the all-unitary lower bound}\label{app:multiscale}
The following compactness argument supplies the lower coefficient
eight. Its quantifier order is essential: the number of scales and
the positive slack are fixed before the dimension limit.

\subsection{A uniform graph bound}
Fix a finite integer $J$ and numbers
$1>\delta_1>\cdots>\delta_J>0$ with
$\delta_{i+1}<\delta_i/2$. Put $\beta_i=2-\delta_i$.
Let $r_i(M)/M^{\beta_i}\to1$ and $r_i c_i^2\to e_i\geq0$.
For arbitrary simple graphs $G_i$ on $M$ vertices with at most $r_i$
edges, put $H_M=\sum_i c_i A(G_i)$, $E=\sum_i e_i$, and
$e_*=\max_i e_i$. Then
\begin{equation}\label{eq:lacunary-graph-bound}
 \limsup_{M\to\infty}\max_{G_1,\ldots,G_J}\lambda_{\max}(H_M)
 \leq\sqrt E+\sqrt{2e_*}.
\end{equation}
We prove uniformity by starting with an arbitrary sequence of graphs
and nonnegative unit maximizing vectors $x$.
Write $p_v=x_v^2$ and $t_v=1+\log_M p_v\leq1$ when $p_v>0$.
Fix $\eta>0$. At level $i$, products $x_vx_w$ below
$M^{-\beta_i/2-\eta}$ contribute at most
$c_i r_i M^{-\beta_i/2-\eta}=O(M^{-\eta})$.
Products above $M^{-\beta_i/2+\eta}$ occur in at most
$M^{\beta_i-2\eta}/2$ unordered pairs, since
$\sum_{v<w}p_vp_w\leq1/2$; Cauchy--Schwarz bounds their weighted
contribution by $O(M^{-\eta})$.
Thus only pairs with $|t_v+t_w-\delta_i|\leq2\eta$ remain.
Coordinates with $t_v<-\eta$ carry total mass at most $M^{-\eta}$.
For edges touching them, Cauchy--Schwarz and
$c_i\sqrt{r_i}=O(1)$ give error $O(M^{-\eta/2})$.

Replace the remaining negative $t_v$ by zero and form the probability
measure $\mu_M=\sum_v p_v\delta_{\max(t_v,0)}$ on $[0,1]$.
Along a subsequence these measures converge weakly to $\mu$, as do
their products. The closed-strip upper bound and Cauchy--Schwarz,
followed by $\eta\downarrow0$, yield
\begin{equation}\label{eq:measure-bound}
 \limsup\tfrac12x^TH_Mx\leq\sum_i\sqrt{e_iQ_i},
 \qquad Q_i=\tfrac12(\mu\times\mu)\{(t,s):t+s=\delta_i\}.
\end{equation}
The same subsequence and measure apply to every level because $J$
is finite. Continuity from above justifies the shrinking strips.
Only atoms contribute to an exact-sum line: its product measure is
the integral of $\mu(\{\delta_i-t\})$, which vanishes away from the
countable set of atoms.

Write $p_t=\mu(\{t\})$. Join distinct atoms when their sum is
$\delta_i$, and allow a loop at $\delta_i/2$. The off-diagonal graph
is a forest. In a finite simple cycle choose an edge of largest
label. Each neighboring edge has a different, smaller label,
hence less than half that label; the two endpoints of the largest
edge would then sum to less than its label, a contradiction.
Consequently
\[
 Q_i=\sum_{t<s,\ t+s=\delta_i}p_tp_s+\tfrac12p_{\delta_i/2}^2.
\]
For $Q_i>0$, set weights proportional to
$\sqrt{p_tp_s}$ on its edges and $p_{\delta_i/2}/\sqrt2$ on its loop,
scaled so their squared sum is $e_i$. These weights express
$\sqrt{e_iQ_i}$ exactly; use zero weights if $Q_i=0$.
The forest is bipartite. Its weighted adjacency operator has norm
at most the square root of the sum of squared edge weights:
after bipartition, its off-diagonal block has operator norm at most
its Hilbert--Schmidt norm. With $z_t=\sqrt{p_t}$ and $\sum_tz_t^2\leq1$,
the total off-diagonal contribution is at most $\sqrt E/2$.
Each loop weight is at most $\sqrt{e_*}$, so all loops contribute
at most $\sqrt{e_*/2}\sum_t p_t\leq\sqrt{e_*/2}$.
The countable case follows by finite truncation; the edge weights
are square-summable. Equation~\eqref{eq:measure-bound} proves
\eqref{eq:lacunary-graph-bound}.

If $a_M\to a$ and $E_f$ is the adjacency matrix of one distinguished
edge, the same argument gives
\begin{equation}\label{eq:isolated-spike}
 \limsup\max_{G_i,f}\lambda_{\max}(H_M+a_ME_f)
 \leq\max\{a,\sqrt E+\sqrt{2e_*}\}.
\end{equation}
Indeed, split off the two endpoints. Their block tends to the
single-edge block of norm $a$, and the cross block has norm at most
$\sqrt{2M}\sum_i c_i\to0$ because every $\beta_i>1$.
The remaining block satisfies the preceding graph bound.

\subsection{Actual rational APPT states}
Fix $h\geq1$, $J=2h^2$, $\delta_i=3^{-i}$, and rational
$0<\varepsilon<1$. Set
\[
 r_i=\lfloor m^{\,2-3^{-i}}\rfloor,\quad
 L_i=\lceil\sqrt{2r_i}\rceil,\quad
 c_i=\frac{2(1-\varepsilon)}{(h+1)L_i},\quad a=2(1-\varepsilon).
\]
The ranks are strictly increasing and below $m(m-1)/2$ eventually.
They have an exact integer definition: $r_i$ is the floor of the
$3^i$th root of $m^{2\cdot3^i-1}$.
Choose nested coordinate projections $P_1\leq\cdots\leq P_J$ of
ranks $r_i$ and a coordinate unit vector $v$ in their first range.
Normalize
\[
 B=I+\sum_i c_iP_i+a|v\rangle\langle v|,\qquad \rho=B/\Tr B.
\]
Here $r_i c_i^2\to e=2(1-\varepsilon)^2/(h+1)^2$,
$E=4h^2(1-\varepsilon)^2/(h+1)^2$, and
$\sqrt E+\sqrt{2e}=2(1-\varepsilon)$.

For every global unitary $U$ and unit Schmidt test vector $x$, a
rank-$r_i$ projection has trace against its partially transposed
projector at least minus the sum of the $r_i$ largest pair products
$x_vx_w$. This follows by evaluating in the witness eigenbasis:
diagonal weights of a rank-$r_i$ projection are in $[0,1]$, with
sum $r_i$. The spike term is bounded below by minus $a$ times the
largest pair product. Thus
\[
 \langle x,(UBU^*)^\Gamma x\rangle
 \geq1-\tfrac12x^T\left(\sum_i c_iA(G_i)+aE_f\right)x
\]
for suitable graphs with at most $r_i$ edges.
Equation~\eqref{eq:isolated-spike}, uniformly over all graphs and
tests, has limit strictly below two. Therefore $\rho$ is APPT for
every sufficiently large $m$ and every $n\geq m$.
No boundary inference is made from a limit equal to two.

Writing $B=I+X$, let $T=\Tr X$ and $Q=\Tr X^2$. Nestedness gives
\[
 T=a+\sum_i c_ir_i,\qquad
 Q=a^2+\sum_i c_i^2r_i+2\sum_{i<k}c_ic_kr_i+2a\sum_i c_i.
\]
As $m\to\infty$, $r_i/r_k\to0$ for $i<k$, so all cross terms vanish;
also $T^2/m^2\to0$ and $T/m^2\to0$. Exactly,
\[
 D^2\bigl(\Tr(\rho^2)-1/D\bigr)
 =\frac{Q-T^2/D}{(1+T/D)^2}.
\]
For $D\geq m^2$ this is at least its value with $D=m^2$,
because $X\geq0$ makes the numerator nonnegative and increasing
and the denominator decreasing as $D$ increases. Its limit there is
$(1-\varepsilon)^2[4+4h^2/(h+1)^2]$.
First choose $h$ and $\varepsilon$ to approach eight, then take $m$
large. The threshold is independent of $n$. This proves the uniform
lower bound eight without interchanging a growing hierarchy and the
dimension limit.

\section{The other lower construction and fixed smaller dimension}\label{app:uniform-lower}
\subsection{A broad plateau with a spike}
For $m\geq4$, set $h=\lceil\sqrt m\rceil$,
\[
 k=\left\lceil\frac{(m-1)n}{2}\right\rceil,\quad
 c=\frac2{m+h},\quad
 a=2-\frac{4(m-2)}{(m+h)(h+3)},\quad q=2-(m-3)c.
\]
Normalize $B=I+cP+(a-c)|v\rangle\langle v|$, with
$\operatorname{rank}P=k$ and $v$ in its range.
The projection and Rayleigh bounds used above reduce every
Schmidt test to the symmetric matrix with diagonal two, off-diagonal
$-a$ on one designated edge and $-c$ on all other edges.
Its endpoint-difference eigenvalue is $2+a$ and its tail-sum-zero
eigenvalue is $2+c$. On the remaining two-dimensional space, positivity
is equivalent to
\[
 2-a\geq0,\qquad q\geq0,\qquad (2-a)q\geq2(m-2)c^2.
\]
The parameters satisfy $q=2(h+3)/(m+h)>0$,
$(2-a)q=2(m-2)c^2$, and
$(a-c)q=[2-(m-1)c](2+c)\geq0$.
Thus $B$ is APPT at every finite pair of dimensions in this range.

Put $y=k-1$, $\theta=y/D$, and $\gamma=n/m$. Its exact centered
purity is $N/Q_0$, where
\[
 N=a^2(1-1/D)+Dc^2\theta(1-\theta)-2a\theta c,\qquad
 Q_0=(1+a/D+\theta c)^2.
\]
Uniformly in $n\geq m$,
$1/2-1/(2m)-1/m^2\leq\theta\leq1/2$. Consequently
\[
 N/Q_0\geq (A_m+B_m\gamma)/Q_m,
\]
for all sufficiently large $m$, where
\begin{align*}
 A_m&=a^2(1-1/m^2)-ac\to4,\\
 B_m&=m^2c^2\left[\tfrac14-(1/(2m)+1/m^2)^2\right]\to1,\\
 Q_m&=(1+a/m^2+c/2)^2\to1.
\end{align*}
Together with Appendix~\ref{app:multiscale}, this supplies the lower
bound in Theorem~\ref{thm:asymptotic}.

\subsection{\texorpdfstring{Fixed $m$ and the all-regime conclusion}{Fixed m and the all-regime conclusion}}
For fixed $m\geq2$ the maximally entangled test gives
$\sum_{i\leq R}\lambda_i\leq\sum_{i\geq D-S+1}\lambda_i$.
Hence $\lambda_R\leq rb$, $r=(m+1)/(m-1)$.
A single-edge arrangement gives $\lambda_1-\lambda_D\leq2b$,
so $\lambda_1\leq3b$. One has $b>0$: if $b=0$, this last
inequality forces every eigenvalue to vanish.
There are $N_0=D-m^2+2$ eigenvalues from index $R$ to index $D-S+1$,
all in $[b,rb]$. Let their sum be $T>0$. The inequality
$x^2\leq(r+1)bx-rb^2$ gives
\[
 \frac{N_0\sum x_i^2}{T^2}
 \leq(r+1)z-rz^2\leq\frac{(r+1)^2}{4r}
 =\frac{m^2}{m^2-1},\qquad z=N_0b/T.
\]
Since $T\leq1$ and the other $m^2-2$ entries are at most $3b\leq3/j$,
\[
 P_{\max}(m,n)\leq\frac{m^2/(m^2-1)}{D-m^2+2}
 +\frac{9(m^2-2)}{(D-S+1)^2}.
\]
For attainment of the limiting coefficient, set $c=2/(m-1)$,
$k=\lceil D/(c+2)\rceil$, and normalize $I+cP$ for a rank-$k$
projection. The negative Schmidt mass is at most $(m-1)/2$,
so this is genuinely APPT. With $\theta=k/D$ its scaled purity
$[1+(2c+c^2)\theta]/(1+c\theta)^2$ has exact deficit
\[
 \frac{(c+2)^2}{4(c+1)}
 -\frac{1+(2c+c^2)\theta}{(1+c\theta)^2}
 =\frac{c^2[1-(c+2)\theta]^2}{4(c+1)(1+c\theta)^2}.
\]
This proves the fixed-$m$ limit and, with the upper bound,
$P_{\max}=m^2/[(m^2-1)D]+O_m(D^{-2})$.
For the single all-regime equivalent, any allegedly failing sequence
has either a fixed-$m$ subsequence or a subsequence with $m\to\infty$.
The first uses this calculation; the second uses
Theorem~\ref{thm:asymptotic} and
\[
 1\leq
 \frac{\max\{8,4+D/(m^2-1)\}}{\max\{8,4+D/m^2\}}
 \leq\frac{m^2}{m^2-1}.
\]
Both contradict a positive limiting error. This proves
Theorem~\ref{thm:allregime}.

\section{The outer-polytope optimization and rectangular equality cases}\label{app:rectangular}
Let $m\geq3$, $n\geq m$, $L=D-m$, and
$u_i=(1/i\ [i\text{ copies}],0\ [D-i\text{ copies}])$.
The ordered simplex has unique gap expansion
$\lambda=\sum_i x_i u_i$, where
$x_i=i(\lambda_i-\lambda_{i+1})$, $\lambda_{D+1}=0$.
The left side minus the right side of~\eqref{eq:outer-single-test}
at $u_i$ is $h_i/i$, with
\[
 h_i=\min(i,R)-\max\{0,i-(D-S)\}.
\]
It is positive for $i<L$, zero at $L$, and negative for $i>L$.
The surviving vertices are $u_j$ for $j\geq L$ and intersections
of the cutting plane with edges $[u_i,u_j]$, $i<L<j$.
For completeness, an on-plane point with three positive gap
coordinates admits a nonzero perturbation preserving the two
affine equations; an off-plane point with two positive coordinates
admits a small perturbation preserving the simplex sum.
Thus there are no other vertices. Every point of this compact
polytope is a convex combination of vertices.

Write $j=L+s$, $1\leq s\leq m$, and $h=h_i>0$.
An edge intersection has spectrum
\[
 \alpha=\frac{h+s}{hj+si}\quad(i\text{ copies}),\qquad
 \beta=\frac h{hj+si}\quad(j-i\text{ copies}),\qquad
 0\quad(D-j\text{ copies}).
\]
These formulas follow by solving trace one and the zero cut value;
the increment has cut value $h(\alpha-\beta)$ and the baseline
has cut value $-s\beta$. Surviving $u_j$ have purity at most $1/L$.

\subsection{The three edge ranges}
For $i\leq R$, $h=i$ and the purity is
\[
 \frac{j+2s+s^2/i}{(j+s)^2}
 \leq f(s):=\frac{L+3s+s^2}{(L+2s)^2}.
\]
It is strict for $i>1$. The derivative
$f'(s)=[(2L-6)s-L]/(L+2s)^3$ changes sign at most once,
from negative to positive. Hence the maximum on $[0,m]$
is $\max\{1/L,Q_{\rm sp}\}$.

For $R\leq i\leq D-S$, $h=R$ and the purity is
\[
 F_i(s)=\frac{R^2L+(R^2+2Ri)s+is^2}{[RL+(R+i)s]^2}.
\]
Its derivative has the sign of
$(2Li-R^2-3Ri-2i^2)s-LR^2$, so no interior maximum is possible.
The endpoints are $1/L$ and
\[
 F_i(m)=\frac{D(m-1)^2+4mi}{[D(m-1)+2i]^2}.
\]
Its derivative in $i$ is
$4m[(m-1)n-2i]/[D(m-1)+2i]^3$.
The real maximum $i_*=(m-1)n/2$ lies in the valid middle interval.
For half-integer $i_*$, with $H=(m^2-1)n$,
\[
 F_{\lceil i_*\rceil}(m)-F_{\lfloor i_*\rfloor}(m)
 =\frac{4m}{(H^2-1)^2}>0.
\]
Thus the unique integer optimum is $t=\lceil i_*\rceil$.
For $D-S<i<L$, $h=L-i$, $\alpha=1/L$, and the purity is
$[i+h^2/(h+s)]/L^2<1/L$.
Finally,
\[
 Q_{\rm sp}-1/L=
 \frac{m[(m-1)D-m(m+3)]}{(D-m)(D+m)^2}\geq0
\]
for $D\geq m^2$ and $m\geq3$.
This proves the full outer maximum $\max\{Q_{\rm sp},Q_t\}$.

\subsection{Threshold and strict equality}
Set $B=m^4-m^2-2m=mn_0(m)$, $u=m^2-1$, and
$C=m^2/(uD)$, the continuous middle maximum. Exact subtraction gives
\[
 C-Q_{\rm sp}=\frac{D^2-BD+m^4}{uD(D+m)^2}.
\]
If $m\leq n<n_0(m)$, the numerator is negative throughout
$m^2\leq D\leq B-m$: it is convex and its endpoint values are
$-m^3(m-2)(m+1)^2$ and $-m^2(m^3-m^2-m-3)$.
The latter bracket is positive for $m=3+y$, $y\geq0$,
since it becomes $y^3+8y^2+20y+12$.
Thus $Q_t\leq C<Q_{\rm sp}$ there.
If $n\geq n_0(m)$, the numerator is at least $m^4$.
When $(m-1)n$ is even there is no rounding loss; otherwise,
with $H=uD/m$,
\[
 Q_t=\frac{m(H+2)}{(H+1)^2},\qquad
 0<C-Q_t=\frac{m}{H(H+1)^2}\leq\frac{m^4}{u^3D^3}.
\]
It follows that
\[
 Q_t-Q_{\rm sp}\geq
 \frac{m^4}{uD(D+m)^2}-\frac{m^4}{u^3D^3}>0,
\]
because $uD>D+m$.
All other vertex ranges are then strictly suboptimal, and the middle
integer optimum is unique. The squared norm is strictly convex:
for distinct $v,w$ and $0<\theta<1$, its Jensen gap equals
$\theta(1-\theta)\|v-w\|^2>0$.
Thus no mixture of distinct vertices can be another maximizer.
The physical projection construction in Theorem~\ref{thm:rectangular}
attains the unique outer vertex. This proves the stated maximum and
the complete unitary-orbit classification in that range.

\section{Certificate file identifiers}\label{app:data}
The following SHA-256 digests identify the exact certificate data. The ancillary manifest also records the checker and this preprint's other source files.
\begingroup\footnotesize
\begin{longtable}{p{0.96\textwidth}}
\toprule
\textbf{Filename and SHA-256}\\\midrule
\endhead
\texttt{qutrit\_cubic\_certificate\_search.json}\\
\nolinkurl{0e5fb10986416ab1a8a5420d92a7a689b5a55d063c9a417d6ce53c2960f1ad77}\\[2pt]
\texttt{qutrit\_D12\_cubic\_certificate.json}\\
\nolinkurl{8ebdf6d6b3712018f91818c6d2bf6336d6a7c6befb34245bf12d719fd3dd31cb}\\[2pt]
\texttt{qutrit\_D15\_fast\_certificate.json}\\
\nolinkurl{c5f688e354e3d2aeabbde72ea343316df64867d60dd095c6f8ee201b466d3df8}\\[2pt]
\texttt{qutrit\_D18\_fast\_certificate.json}\\
\nolinkurl{1312bfc0488d8f7ca8afff43007d93d62d81dc9a0cd664f57196b4622d328811}\\[2pt]
\texttt{qutrit\_D21\_fast\_certificate.json}\\
\nolinkurl{34931c9638725825a0efe8bd80d4a1cb15e53de6d1b433b8b77809f09597ea46}\\[2pt]
\texttt{qutrit\_D24\_fast\_certificate.json}\\
\nolinkurl{b19f93c12f369530ff436f095700efa631b029a8964a29e32fb3ea125d22d65e}\\[2pt]
\texttt{qutrit\_uniform\_parameter\_search.json}\\
\nolinkurl{d381b47b3bcb67ddc457f7c866b5a6e97757e81cb9d7a84fb832a7b397a826a3}\\
\bottomrule
\end{longtable}\endgroup

\begin{thebibliography}{9}\small
\bibitem{Hildebrand} R. Hildebrand, Positive partial transpose from spectra, \emph{Physical Review A} \textbf{76} (2007), 052325. \href{https://doi.org/10.1103/PhysRevA.76.052325}{doi:10.1103/PhysRevA.76.052325}. Preprint: \href{https://arxiv.org/abs/quant-ph/0502170}{arXiv:quant-ph/0502170}.
\bibitem{DungKhoi} H. P. Dung and V. T. Khoi, On the maximum purity of absolutely separable bipartite states, preprint (2025), \href{https://arxiv.org/abs/2510.19508v1}{arXiv:2510.19508v1}.
\bibitem{AKW} J. Ahiable, N. B. T. Kothakonda and A. Winter, The geometry of absolute separability and other convex matrix properties from spectrum, preprint (2026), \href{https://arxiv.org/abs/2608.03390v2}{arXiv:2608.03390v2}. In particular, Theorem 6.3, Corollary 6.4 and Conjecture 6.7.
\bibitem{Tran} A. T. Tran, Spectral optimization for absolutely PPT states: purity, entropy, and volume decay, preprint (2026), \href{https://arxiv.org/abs/2609.18568}{arXiv:2609.18568}.
\bibitem{Kondra} T. V. Kondra, P. Barrios Hita, J. Neumann, H. Kampermann and D. Bru{\ss}, Fundamental limitations on entanglement extraction from purity, preprint (2026), \href{https://arxiv.org/abs/2605.29197v1}{arXiv:2605.29197v1}. Supplemental Material, Lemmas 1--2.
\bibitem{ProofRelease} Y. Zhang, Qutrit--qudit APPT purity: complete all-$n$ Lean proof v1, immutable software and proof-data release (2026), \texttt{appt-qutrit-purity-complete-v1}. \href{https://github.com/mxym/math/releases/tag/appt-qutrit-purity-complete-v1}{GitHub immutable proof release}; publication commit \proofcommit; \href{https://github.com/mxym/math/actions/runs/37968137659}{successful audit 37968137659}.
\end{thebibliography}
\end{document}
